\chapter{Uitdrukkingen met letters}\label{ch:g7:literal}

Een letter in een berekening staat voor een getal dat we nog niet kennen, of voor
\emph{alle} getallen tegelijk: $P = 4z$ schrijven geeft de \omterm{def:g6:measure:perimeter}{omtrek} van elk vierkant
in één formule. Dit hoofdstuk leert de schrijfafspraken, het invullen van waarden
en het vereenvoudigen van eenvoudige uitdrukkingen --- het begin van de algebra,
die in \cref{ch:g8:equations} verdergaat.

\section{Schrijven met letters}

\begin{definition}[Uitdrukking met letters]\label{def:g7:literal:def}
Een \emph{uitdrukking met letters}\index{uitdrukking met letters} is een
berekening met een of meer letters, die elk voor een getal staan. De
schrijfafspraken:
\begin{itemize}
  \item het teken $\times$ valt weg voor een letter of een haakje:
        $3 \times a = 3a$, \ $a \times b = ab$, \
        $2 \times (x + 5) = 2(x + 5)$;
  \item maar het blijft staan tussen twee getallen: $3 \times 5$ mag je niet als
        $35$ schrijven!
  \item $a \times a$ schrijf je als $a^2$ (``$a$ kwadraat''), en
        $a \times a \times a = a^3$;
  \item $1a$ schrijf je als $a$, en $0a$ is $0$.
\end{itemize}
\end{definition}

\begin{example}\label{ex:g7:literal:write}
Vereenvoudig de schrijfwijze:
\[
5 \times x + 3 \times y = 5x + 3y,
\qquad
a \times 7 = 7a \ \text{(getal vooraan)},
\qquad
x \times x \times 4 = 4x^2 .
\]
\end{example}

\begin{example}[Formules zijn uitdrukkingen met letters]\label{ex:g7:literal:formulas}
Voor een rechthoek met lengte $l$ en breedte $b$: \omterm{def:g6:measure:perimeter}{omtrek} $P = 2(l + b)$,
\omterm{def:g6:measure:area}{oppervlakte} $A = lb$. Eén formule vat de oneindig veel rechthoeken samen --- en
dat is precies waar letters voor zijn.
\end{example}

\begin{omfigure}
\begin{tikzpicture}[scale=0.8]
  \draw[thick] (0,0) rectangle (4.4,2.0);
  \node[below, font=\small] at (2.2,0) {$l$};
  \node[left, font=\small] at (0,1.0) {$b$};
  \node[font=\small] at (2.2,1.0) {$A = lb$};
  \node[above, font=\small, omDef] at (2.2,2.05) {$P = 2(l + b)$};
\end{tikzpicture}

{\small Twee formules die voor elke rechthoek geldig zijn: vul de waarden van $l$
en $b$ in om getallen te krijgen.}
\end{omfigure}

\section{Invullen}

\begin{method}[Een uitdrukking uitrekenen]\label{met:g7:literal:substitute}
Om een uitdrukking voor gegeven waarden van haar letters uit te rekenen:
\begin{enumerate}
  \item schrijf de uitdrukking met de verborgen $\times$-tekens erbij:
        $3a + a^2$ is $3 \times a + a \times a$;
  \item vervang elke letter door haar waarde, met haakjes eromheen als dat nodig
        is;
  \item reken uit, met de voorrangsregels van \cref{ch:g7:priorities}.
\end{enumerate}
\end{method}

\begin{example}\label{ex:g7:literal:substitute}
Reken $E = 3x + x^2$ uit voor $x = 5$:
\[
E = 3 \times 5 + 5 \times 5 = 15 + 25 = 40 .
\]
Reken $F = 2(a + 3b)$ uit voor $a = 4$ en $b = 0.5$:
\[
F = 2 \times (4 + 3 \times 0.5) = 2 \times (4 + 1.5) = 2 \times 5.5
= 11 .
\]
\end{example}

\section{Uitdrukkingen herleiden}

\begin{proposition}[Gelijksoortige termen samennemen]\label{prop:g7:literal:reduce}
Termen met hetzelfde letterdeel kun je samennemen door hun getallen op te tellen
(dat is de distributiviteit, \cref{thm:g7:priorities:distributivity},
achterstevoren gelezen):
\[
5x + 3x = (5 + 3)x = 8x,
\qquad
7a - 2a = 5a .
\]
Termen met verschillende letterdelen (zoals $3x$ en $5y$, of $x$ en $x^2$) kun je
niet samennemen.
\end{proposition}

\begin{example}\label{ex:g7:literal:reduce}
Herleid $E = 4x + 7 + 3x + 2$ door de termen te sorteren:
\[
E = (4x + 3x) + (7 + 2) = 7x + 9 .
\]
Pas op: $7x + 9$ is \emph{niet} $16x$! Test met $x = 1$: $4 + 7 + 3 + 2 = 16$ en
$7 \times 1 + 9 = 16$ komen overeen, maar $16x$ zou ook $16$ zijn --- test dus met
$x = 2$: het origineel geeft $8 + 7 + 6 + 2 = 23$, en $7 \times 2 + 9 = 23$,
terwijl $16 \times 2 = 32$. De herleiding $7x + 9$ is juist, $16x$ is fout.
\end{example}

\begin{method}[Een gelijkheid testen]\label{met:g7:literal:test}
Om na te gaan of twee uitdrukkingen voor elke waarde gelijk kunnen zijn:
\begin{enumerate}
  \item vul in beide dezelfde waarde in, minstens twee keer (bijvoorbeeld
        $x = 2$ en $x = 10$; vermijd $0$ en $1$, die verbergen te veel
        verschillen);
  \item verschillen de uitkomsten ook maar één keer, dan zijn de uitdrukkingen
        \emph{niet} gelijk --- één tegenvoorbeeld velt een algemene bewering;
  \item komen de uitkomsten overeen, dan is de gelijkheid alleen
        \emph{aannemelijk}: een bewijs vraagt algebra (uitwerken, herleiden).
\end{enumerate}
\end{method}

\begin{example}\label{ex:g7:literal:test}
Is $2(x + 3)$ gelijk aan $2x + 6$ voor elke $x$? Uitwerken bewijst het:
$2(x+3) = 2x + 2 \times 3 = 2x + 6$. Is $(x + 1)^2$ gelijk aan $x^2 + 1$? Test
$x = 2$: $(2+1)^2 = 9$ maar $2^2 + 1 = 5$. Nee --- en één tegenvoorbeeld is een
volledig bewijs van ``nee''.
\end{example}

\section{Uitdrukkingen opstellen}

\begin{example}[Van situatie naar uitdrukking]\label{ex:g7:literal:produce}
Een taxi rekent $4$ euro plus $2$ euro per kilometer. Voor een rit van $k$
kilometer is de prijs
\[
p = 4 + 2k \quad\text{euro}.
\]
Voor $k = 7$: $p = 4 + 14 = 18$ euro. De uitdrukking \emph{is} het algemene
antwoord; invullen levert elk bijzonder antwoord op.
\end{example}

\begin{example}[Getaltrucs]\label{ex:g7:literal:trick}
``Denk aan een getal, verdubbel het, tel er $10$ bij op, deel door $2$, en haal je
eigen getal eraf.'' Noem het getal $n$ en volg de stappen:
\[
n \ \to\ 2n \ \to\ 2n + 10 \ \to\ \frac{2n + 10}{2} = n + 5 \ \to\
n + 5 - n = 5 .
\]
Iedereen krijgt $5$, wat $n$ ook was --- de letters verklaren de goocheltruc.
\end{example}

\section{Oefeningen}

\begin{exercise}[$\star$]\label{exo:g7:literal:1}
Vereenvoudig de schrijfwijze:
\[
7 \times x, \qquad
y \times 3, \qquad
a \times b \times 5, \qquad
x \times x, \qquad
2 \times (a + 4).
\]
\end{exercise}

\begin{exercise}[$\star$]\label{exo:g7:literal:2}
Zet de $\times$-tekens terug:
\[
5ab, \qquad
3(x + 2), \qquad
x^2 + 4x, \qquad
2a^3 .
\]
\end{exercise}

\begin{exercise}[$\star$]\label{exo:g7:literal:3}
Reken uit voor $x = 3$: \ $5x + 1$; \quad $x^2$; \quad $2x^2 - x$; \quad
$10 - 2x$.
\end{exercise}

\begin{exercise}[$\star$]\label{exo:g7:literal:4}
Reken $3a + 2b$ en $a b + 5$ uit voor $a = 4$ en $b = 1.5$.
\end{exercise}

\begin{exercise}[$\star$]\label{exo:g7:literal:5}
Herleid:
\[
8x + 5x, \qquad
9a - 4a + a, \qquad
6y + 2 + 3y + 7, \qquad
5t + 3 - 2t - 3 .
\]
\end{exercise}

\begin{exercise}[$\star$]\label{exo:g7:literal:6}
Geef een formule voor: de \omterm{def:g6:measure:perimeter}{omtrek} van een vierkant met zijde $z$; de \omterm{def:g6:measure:perimeter}{omtrek} van een
\omterm{def:g6:shapes:triangles}{gelijkbenige driehoek} met basis $b$ en gelijke zijden $s$; de prijs van $n$
broden van $1.20$ euro per stuk.
\end{exercise}

\begin{exercise}[$\star$]\label{exo:g7:literal:7}
Koppel elke uitdrukking aan een beschrijving: $2n$; $n + 2$; $n^2$; $\frac n2$.
Beschrijvingen: ``de \omterm{def:g3:division:half}{helft} van $n$''; ``het dubbele van $n$''; ``het kwadraat van
$n$''; ``twee meer dan $n$''.
\end{exercise}

\begin{exercise}[$\star\star$]\label{exo:g7:literal:8}
Beslis met \cref{met:g7:literal:test} (met een bewijs of een tegenvoorbeeld) of
deze gelijkheden voor elke $x$ gelden:
\[
3(x + 4) = 3x + 12; \qquad
2x + 5 = 7x; \qquad
x + x = x^2 .
\]
\end{exercise}

\begin{exercise}[$\star\star$]\label{exo:g7:literal:9}
Een rechthoek heeft breedte $b$ en zijn lengte is $3$ cm meer dan zijn breedte.
\begin{enumerate}
  \item Druk de lengte uit in $b$, en daarna de \omterm{def:g6:measure:perimeter}{omtrek}, en herleid.
  \item Reken de \omterm{def:g6:measure:perimeter}{omtrek} uit voor $b = 5$ cm, rechtstreeks en met je herleide
        formule. Hetzelfde antwoord?
\end{enumerate}
\end{exercise}

\begin{exercise}[$\star\star$]\label{exo:g7:literal:10}
Volg de truc van \cref{ex:g7:literal:trick} met letters: ``denk aan een getal, tel
er $6$ bij op, vermenigvuldig met $3$, haal er $18$ af, en deel door je startgetal''
(neem aan dat dat niet \omterm{def:g1:counting:zero}{nul} is). Wat krijgt iedereen?
\end{exercise}

\begin{exercise}[$\star\star\star$]\label{exo:g7:literal:11}
De figuur hieronder is een vierkant met zijde $a$ waaruit in een hoek een vierkant
met zijde $b$ is weggehaald ($b$ kleiner dan $a$). Schrijf twee verschillende
uitdrukkingen voor haar \omterm{def:g6:measure:area}{oppervlakte} op (één als \omterm{ex:g1:subtraction:difference}{verschil}, één door de L-vorm in
twee rechthoeken te knippen), en ga na dat ze overeenkomen voor $a = 5$ en
$b = 2$.
\end{exercise}

\section{Opgave: het luciferorakel}

\begin{problem}[{Weekendopgave --- groeiende patronen en hun $n$-de term: drie
formules voor één patroon, en het beroemde patroon dat bij de zesde stap
bezwijkt}]
\label{pb:g7:literal:1}
Leg luciferstokjes in een rij vierkanten: één vierkant, dan twee, dan drie\,\dots\
Hoeveel stokjes heeft de honderdste figuur nodig? Ze één voor één tellen zou je
hele weekend kosten; een \emph{formule in $n$} antwoordt in één keer voor elke
figuur --- precies waar letters voor zijn
(\cref{ex:g7:literal:formulas}). Maar pas op: deze opgave eindigt met een
gevierd patroon dat vijf stappen eerlijk speelt en daarna iedereen verraadt die
erop vertrouwde.

\medskip
\noindent\textbf{Deel I --- Een rij vierkanten.}
Figuur $1$ is één vierkant van stokjes ($4$ stokjes); elke volgende figuur voegt
er één vierkant aan toe dat een zijde met het vorige deelt.
\begin{enumerate}
  \item Teken de figuren $1$ tot $4$ en tel hun stokjes.
  \item Leg uit waarom elk nieuw vierkant precies $3$ extra stokjes kost, en
        voorspel daarmee het aantal voor figuur $10$ zonder formule.
  \item Verantwoord de formule: figuur $n$ heeft $3n + 1$ stokjes nodig. (Waar komt
        die eenzame ``$+1$'' vandaan?) Controleer haar met vraag 1.
  \item Amir denkt anders: ``$4$ stokjes voor het eerste vierkant, en dan $3$ voor
        elk van de $n - 1$ andere'', dus $4 + 3(n - 1)$. Werk zijn uitdrukking uit
        en herleid haar (\cref{prop:g7:literal:reduce}): is hij het met vraag 3
        eens?
  \item Lena telt op een derde manier: ``$2n$ horizontale stokjes en $n + 1$
        verticale.'' Schrijf haar formule op, herleid haar, en besluit. Antwoord
        daarna op de openingsvraag: hoeveel stokjes voor figuur $100$?
\end{enumerate}

\medskip
\noindent\textbf{Deel II --- Driehoeken, tafels en achterstevoren rekenen.}
\begin{enumerate}[resume]
  \item Een rij driehoeken van stokjes (elke nieuwe driehoek leunt tegen de
        vorige): tel de stokjes voor $1$, $2$ en $3$ driehoeken, en verantwoord
        daarna de formule $2n + 1$ voor $n$ driehoeken.
  \item Achterstevoren: hoeveel driehoeken kan de rij hebben met precies $49$
        stokjes? (Maak de formule stap voor stap ongedaan.)
  \item Aan vierkante banketttafels zit één gast per vrije zijde. Leg uit waarom de
        formule voor $n$ tafels in één lange rij $2n + 2$ plaatsen geeft, en
        controleer haar voor $1$, $2$ en $3$ tafels.
  \item Op een feest worden $30$ gasten verwacht. Hoeveel tafels heeft één rij
        nodig? De cateraar zet dezelfde tafels in plaats daarvan in \emph{twee}
        rijen van zeven: hoeveel gasten passen er nu? Verklaar de winst met de
        formule.
  \item Stippenpatroon: figuur $n$ is een L van stippen, $n$ stippen hoog en $n$
        stippen breed (de hoekstip één keer geteld). Tel de stippen van de figuren
        $1$, $2$ en $3$, zoek de formule, en noem de getallen die ze oplevert. (Op
        elkaar gelegd betegelen deze L's een vierkant --- een verhaal dat in
        \cref{pb:g8:equations:1} verdergaat.)
\end{enumerate}

\medskip
\noindent\textbf{Deel III --- Vertrouwen, testen, bewijzen.}
\begin{enumerate}[resume]
  \item Een klasgenoot beweert dat figuur $n$ van deel I $3(n + 2) - 5$ stokjes
        nodig heeft. Beslis over die bewering door uit te werken en te herleiden
        --- testen is niet nodig.
  \item ``Kwadrateren maakt getallen groter, dus $n^2 \geq n$ voor elk getal
        $n$.'' Test de bewering voor $n = 2$, $3$ en $10$; zoek daarna een getal
        waarvoor ze faalt. Wat hebben de drie geslaagde tests bewezen, en wat
        bewijst die ene mislukking (\cref{met:g7:literal:test})?
  \item Zet punten op een cirkel en teken \emph{elk} \omterm{def:g6:lines:objects}{lijnstuk} tussen die punten;
        tel daarna de gebieden binnen de cirkel. Doe het voor $2$, $3$, $4$ en $5$
        punten (zet de punten ongelijkmatig, zodat er nooit drie lijnstukken door
        één binnenpunt gaan). Noteer de aantallen. Welke formule fluistert het
        patroon je in?
  \item Nu $6$ punten, groot en ongelijkmatig getekend, en heel zorgvuldig geteld
        (er zijn veel kleine gebieden). Hoeveel gebieden vind je --- en wat is er
        met het vermoeden gebeurd?
  \item De moraal, in twee zinnen: wat zou er nodig zijn geweest om het
        verdubbelvermoeden te \emph{bewijzen}, en waarom hebben de
        luciferformules van deel I en II nooit in zo'n gevaar verkeerd? Vier het
        met een bewezen voorspelling: hoeveel stokjes voor de rij van $2\,026$
        vierkanten?
\end{enumerate}
\end{problem}
